\begin{abstract}
Optical music recognition (OMR) systems achieve strong performance on in-domain benchmarks, yet whether such results translate into plug-and-play recognition on unseen score domains remains underexplored. We present RIGOR-OMR (Relation-aware Interpretable Geometry for Out-of-domain Recognition), a modular framework that leverages staff-normalized geometry to guide visual symbol detection and explicit relation reasoning, while reconstructing structured outputs without target-domain training or dataset-specific adaptation. To compare systems with heterogeneous outputs, we introduce a representation-agnostic pitch--duration event evaluation protocol. On Polish Scores, RIGOR-OMR achieves 17.083 Event-F1, representing relative improvements of 357.8\% over SMT and 1182.4\% over Zeus. On Debussy, RIGOR-OMR significantly outperforms SMT while performing comparably to Zeus. Ablation studies confirm that explicit relation reasoning is the primary source of these improvements. Extensive experiments across Polish Scores, Debussy, GrandStaff, and OLiMPiC establish RIGOR-OMR as an effective approach to zero-shot music-content recognition under substantial domain shift.
\end{abstract}
